\documentclass[10pt,a4paper]{amsart}
\usepackage[margin=19mm]{geometry}
\usepackage{amsmath,amssymb,mathtools}
\usepackage[scr=rsfs,cal=euler]{mathalfa}
\usepackage{xfrac}
\usepackage[unicode,hidelinks,bookmarks=false]{hyperref}
\newcommand{\ie}{i.e.\ }
\newcommand{\cf}{cf.\ }
\newcommand{\resp}{respectively\ }
\newcommand{\Subsection}{\subsection}
\providecommand{\nobreakdash}{\nobreak\hbox{-}}
\setlength{\emergencystretch}{2em}
\multlinegap0pt
\usepackage{amssymb}
\usepackage{mathtools}
\usepackage[shortlabels]{enumitem}
\usepackage{float}
\usepackage{tikz}
\usepackage{xcolor}
\usepackage[normalem]{ulem}
\usepackage{cancel}
\newtheorem{lemma}{Lemma}
\newcommand{\R}{\mathbb R}
\newcommand{\bq}{\begin{equation}}
\newcommand{\dd}{\mathrm d}
\newcommand{\dpp}{\mathrm{d}p}
\newcommand{\ds}{\mathrm{d}s}
\newcommand{\dt}{\mathrm{d}t}
\newcommand{\dx}{\mathrm{d}x}
\newcommand{\dy}{\mathrm{d}y}
\newcommand{\eq}{\end{equation}}
\newcommand{\lt}{\left}
\newcommand{\one}{\mathbf 1}
\newcommand{\rt}{\right}
\title{Global mild solutions and the semiclassical limit for the Fermi--Dirac BGK model}
\author{Young-Pil Choi, Byung-Hoon Hwang, Sihyun Song}
\date{Source: arXiv:2609.08142v1 (CC BY 4.0)}
\begin{document}
\maketitle

\noindent\emph{Source and licence.} Global mild solutions and the semiclassical limit for the Fermi--Dirac BGK model. Version arXiv:2609.08142v1 (CC BY 4.0), distributed by arXiv under the Creative Commons Attribution 4.0 International licence, http://creativecommons.org/licenses/by/4.0/, as recorded in the arXiv licence metadata for this exact version and shown on the abstract page https://arxiv.org/abs/2609.08142v1. Full licence notice: \texttt{licenses/arxiv-2609.08142-CC-BY-4.0.txt}.\par\smallskip

\noindent\textit{Editorial scope.} Counted unit: the article's lemma lem:3-mom (source0.tex lines 1382--1392). The article's complete original proof of it is delivered with it (source0.tex lines 1394--1476, one \texttt{proof} environment of 83 lines). No mathematics is written, repaired or re-derived: the statement and the proof are the article's own lines, in the article's own notation, and no retained line is substituted or re-flowed.  No further source-local context is needed: every cross-reference of every retained line resolves inside the two delivered spans.  Nothing was omitted from any retained span, so the package carries no omission record.  No retained line contains a citation, so the wrapper carries no bibliography and no bibliography entry.  No retained line contains a figure environment, an image-inclusion command or drawing code, so no image is delivered and the package declares no asset dependency.  Editorial formatting only, in addition: an AMS \texttt{amsart} preamble that restates the article's own declarations and macros used by the retained lines, namely the environments \texttt{lemma} on separate counters, together with the standard packages those lines need.  The retained environments print in this document as Lemma 1 in order of appearance, whereas the article prints the same items with the numbering of its own full document.  The complete native source of the article as distributed in the arXiv e-print for this exact version is bundled unchanged as \texttt{sources/209698/source0.tex}.
\begin{lemma}\label{lem:3-mom}
Let $T,R>0$, and suppose
\[
 f(t,x,p)  =e^{-t}f_0(x-tp,p) +\int_0^te^{-(t-s)}G(s,x-(t-s)p,p)\,\ds,
\]
where $f_0,G\geq0$.  Then
\begin{align*}
 &\int_0^T\int_{B_R}\int_{\R^3}|p|^3f(t,x,p)\,\dpp\dx\dt \\
 &\quad\leq 2R\lt(  \iint_{\R^3\times\R^3}|p|^2f_0(x,p)\,\dx\dpp  +\int_0^T\iint_{\R^3\times\R^3}|p|^2G(s,x,p)\,\dx\dpp\ds  \rt).
\end{align*}
\end{lemma}

\begin{proof}
We first record a geometric estimate for free trajectories. For $y,p\in\R^3$, set
\[
 A_{R,T}(y,p)  :=  \int_0^T \one_{B_R}(y+tp)\,\dt  =  \lt|\{t\in[0,T]:y+tp\in B_R\}\rt|.
\]
Thus $A_{R,T}(y,p)$ is the amount of time that the free trajectory $t\mapsto y+tp$ spends in $B_R$ during $[0,T]$.

Suppose first that $p\neq0$, and set
\[
 e:=\frac{p}{|p|},  \quad  y_\perp:=y-(y\cdot e)e.
\]
Then, we have $ y+tp  =  y_\perp+(y\cdot e+t|p|)e$, and consequently
\[
 |y+tp|^2  =  |y_\perp|^2+(y\cdot e+t|p|)^2
\]
due to $y_\perp\cdot e=0$. If $|y_\perp|\geq R$, then the line $\{y+tp:t\in\R\}$ does not enter $B_R$, up to the irrelevant tangential case, and hence
\[
 A_{R,T}(y,p)=0.
\]
If $|y_\perp|<R$, set $\rho:=\sqrt{R^2-|y_\perp|^2}$. The condition $y+tp\in B_R$ is then equivalent to $ |y\cdot e+t|p||<\rho$, that is,
\[
 \frac{-\rho-y\cdot e}{|p|}  <  t  <  \frac{\rho-y\cdot e}{|p|}.
\]
Thus the set of all $t\in\R$ for which $y+tp\in B_R$ is an interval of length
\[
 \frac{2\rho}{|p|}  =  \frac{2\sqrt{R^2-|y_\perp|^2}}{|p|}.
\]
Restricting this interval to $[0,T]$, we obtain
\bq \label{eq:crossing-time}
 A_{R,T}(y,p)  \leq  \min\lt\{  T,\,  \frac{2\sqrt{(R^2-|y_\perp|^2)_+}}{|p|}  \rt\}  \leq  \min\lt\{T,\frac{2R}{|p|}\rt\}.
\eq
In particular,
\bq \label{eq:crossing-moment}
 |p|^3A_{R,T}(y,p)  \leq  \min\{T|p|^3,2R|p|^2\}  \leq  2R|p|^2.
\eq
For $p=0$, we simply have
\[
 A_{R,T}(y,0)=T\one_{B_R}(y),
\]
hence \eqref{eq:crossing-moment} remains valid.

We now estimate separately the two terms in the mild formula. For the contribution of the initial datum, Tonelli's theorem, the change of variables $y=x-tp$, and $e^{-t}\leq1$ give
\begin{align*}
 \int_0^T\int_{B_R}\int_{\R^3}  |p|^3e^{-t}f_0(x-tp,p)\,\dpp\dx\dt  &=  \iint_{\R^3\times\R^3}  |p|^3f_0(y,p)  \lt(  \int_0^T e^{-t}\one_{B_R}(y+tp)\,\dt  \rt)  \,\dy\dpp  \\
 & \leq  \iint_{\R^3\times\R^3}  |p|^3f_0(y,p)A_{R,T}(y,p)\,\dy\dpp  \\
 & \leq  2R  \iint_{\R^3\times\R^3}  |p|^2f_0(y,p)\,\dy\dpp,
\end{align*}
where the last inequality follows from \eqref{eq:crossing-moment}.  Notice that this bound is uniform in $T$: a fixed free trajectory can spend at most $2R/|p|$ units of time in $B_R$, independently of the length of the observation interval.

We next consider the source term.  By Tonelli's theorem,
\begin{align*}
 I_G  &:=  \int_0^T\int_{B_R}\int_{\R^3}|p|^3  \int_0^t e^{-(t-s)}  G(s,x-(t-s)p,p)\,\ds  \,\dpp\dx\dt  \\
 &=  \int_0^T\int_s^T\int_{B_R}\int_{\R^3}  |p|^3e^{-(t-s)}  G(s,x-(t-s)p,p)  \,\dpp\dx\dt\ds.
\end{align*}
For fixed $s$, set
\[
 \tau:=t-s,  \quad  y:=x-\tau p.
\]
Then $0\leq\tau\leq T-s$, while $x\in B_R$ is equivalent to $y+\tau p\in B_R$.
Thus,
\begin{align*}
 I_G  &=  \int_0^T  \iint_{\R^3\times\R^3}  |p|^3G(s,y,p)  \lt(  \int_0^{T-s}  e^{-\tau}\one_{B_R}(y+\tau p)\,\dd\tau  \rt)  \,\dy\dpp\ds  \\
 &\leq  \int_0^T  \iint_{\R^3\times\R^3}  |p|^3G(s,y,p)  A_{R,T-s}(y,p)  \,\dy\dpp\ds  \\
 &\leq  2R  \int_0^T  \iint_{\R^3\times\R^3}  |p|^2G(s,y,p)\,\dy\dpp\ds,
\end{align*}
where we used \eqref{eq:crossing-moment}, with $T-s$ in place of $T$, in the last step.

Combining the two estimates gives
\[
\int_0^T\int_{B_R}\int_{\R^3}  |p|^3f(t,x,p)\,\dpp\dx\dt  \leq  2R\lt(  \iint_{\R^3\times\R^3}  |p|^2f_0(x,p)\,\dx\dpp  +  \int_0^T  \iint_{\R^3\times\R^3}  |p|^2G(s,x,p)\,\dx\dpp\ds  \rt),
\]
which is the desired estimate.

We emphasize that the geometric estimate \eqref{eq:crossing-moment} itself is uniform in $T$. The dependence on $T$ in the final bound enters only through the time integral of the source term.  For example, if
\[
 \sup_{0\leq s\leq T}  \iint_{\R^3\times\R^3}  |p|^2G(s,x,p)\,\dx\dpp  \leq C_T,
\]
then
\[
 \int_0^T\int_{B_R}\int_{\R^3}  |p|^3f(t,x,p)\,\dpp\dx\dt  \leq  2R\lt(  \iint_{\R^3\times\R^3}|p|^2f_0\,\dx\dpp  +TC_T  \rt).
\]
This completes the proof.
\end{proof}

\end{document}
